Scale temporarily to side one, with $A=(0,0),B=(1,0),C=(1,1),D=(0,1),P=(0,t)$. Parameterize line $BP$ as $(1-u,ut)$. Its intersections with $CE,CF$ have parameters $u_X=\dfrac{1}{t+2}$ and $u_Y=\dfrac{1}{t+4}$. Hence $XY=2YB$ is $\left|\dfrac{2}{t+2}\right|=2$, giving $t=-1,-3$. Also $G=(\dfrac{2}{3},\dfrac{1}{3})$ and $H=(\dfrac{4}{5},\dfrac{1}{5})$. Their distance is $\dfrac{2\sqrt2}{15}$, and the distances of the two possible $P$ to line $BD$ are $\sqrt2,2\sqrt2$. The areas are $\dfrac{2}{15},\dfrac{4}{15}$ in unit scale. Multiplying by $45^2$ gives $\boxed{270\text{ or }540}$. 
\begin{asy}
texpreamble("\usepackage{amsmath,amssymb}");
size(220);
defaultpen(fontsize(9pt));
for(int k=0;k<2;++k){pair v=(3*k,0),A=v,B=v+(1,0),C=v+(1,1),D=v+(0,1),P=v+(0,k==0?-1:-3),G=v+(2/3,1/3),H=v+(4/5,1/5);draw(A--B--C--D--cycle);draw(B--D);draw(P--G--H--cycle);label("$P$",P,(0,-1));label("$G$",G,(-1,0));label("$H$",H,(1,0));label("$A$",A,(-1,0));label("$B$",B,(1,0));label("$C$",C,(1,1));label("$D$",D,(-1,1));}
\end{asy}
